\section{Large Language Monkeys: Scaling Law del muestreo repetido}

\subsection{El sistema de dos etapas: generación y verificación}

Repeated sampling primero muestrea de manera independiente múltiples candidatos a partir del mismo modelo, y después los entrega a un verifier para que elija. El código y las demostraciones formales se adaptan especialmente bien a este método, porque las pruebas unitarias, los compiladores y el proof checker pueden ofrecer retroalimentación automática casi determinista.

\videofigure{figures/generate-verify.jpg}{La generación de múltiples candidatos y su verificación constituyen el flujo básico del inference scaling.}{00:01:18--00:02:22}

Supongamos que la probabilidad de éxito de un solo muestreo para el problema $i$ es $p_i$; la oracle coverage al muestrear $k$ veces de forma independiente es:

$$
C_i(k)=1-(1-p_i)^k.
$$

\begin{itemize}
    \item $p_i$: la probabilidad de éxito de un solo intento del modelo en el problema $i$;
    \item $k$: el número de muestreos;
    \item $C_i(k)$: la probabilidad de que el conjunto de candidatos contenga al menos una solución correcta.
\end{itemize}

Esto muestra que cada problema en sí mismo presenta una saturación exponencial: el muestreo inicial aporta el mayor beneficio, y después la ganancia marginal disminuye gradualmente.

\subsection{Por qué un modelo pequeño puede alcanzar a un modelo fuerte mediante el muestreo}

Mientras $p_i>0$, aumentar $k$ puede elevar la coverage. En los experimentos, un modelo más pequeño reduce de manera notable, mediante un gran número de muestreos, la brecha respecto de un modelo fuerte con inferencia única en matemáticas, código y demostraciones formales.

\videofigure{figures/repeated-sampling.jpg}{El muestreo repetido amplifica de manera notable la capacidad del modelo en múltiples reasoning benchmarks.}{00:02:22--00:04:06}

\begin{knowledgebox}{¿La capacidad se "amplifica" o se "descubre"?}
Repeated sampling no modifica los parámetros ni hace que el modelo aprenda un algoritmo nuevo. Lo que hace es aumentar, dentro de la distribución de salidas que el modelo ya posee, la probabilidad de obtener una trayectoria correcta de baja probabilidad. Por eso es más preciso decir que \textbf{descubre y aprovecha una capacidad latente}, y no que crea una capacidad totalmente nueva.
\end{knowledgebox}

\subsection{De la curva exponencial de un solo problema a la ley de potencias promedio del dataset}

El ajuste empírico que ofrece el curso es la exponentiated power law:

$$
\bar C(k)\approx 1-\exp(-a k^b),\qquad a>0,\;0<b<1.
$$

\begin{itemize}
    \item $\bar C(k)$: la coverage promedio sobre el dataset;
    \item $a$: un parámetro de escala relacionado con el modelo y la dificultad general de la tarea;
    \item $b$: el exponente que controla la velocidad del scaling;
    \item $k$: el presupuesto de muestreo en inferencia.
\end{itemize}

\videofigure{figures/inference-scaling-law.jpg}{La coverage promedio sigue una exponentiated power law predecible conforme aumenta el número de muestreos.}{00:05:18--00:07:06}

Un solo problema es $1-(1-p_i)^k$; ¿por qué al promediar sobre los problemas surge una ley de potencias? La clave está en la distribución de $p_i$: en el dataset existe una gran cantidad de problemas difíciles que el modelo solo resuelve ocasionalmente, lo que forma una cola larga cercana a cero.

\videofigure{figures/scaling-across-models.jpg}{El inference scaling se observa tanto en modelos pequeños como en modelos grandes.}{00:07:07--00:09:40}

\subsection{La cola larga es una condición suficiente para la ley de potencias promedio}

Si la probabilidad de éxito de los problemas tiene, cerca de $p\rightarrow0$, una densidad

$$
f(p)\propto p^{\alpha-1},
$$

entonces la tasa de fallo promedio del dataset se aproxima a:

$$
1-\bar C(k)=\mathbb{E}_{p}[(1-p)^k]
\approx\int_{0}^{\infty} e^{-kp}p^{\alpha-1}\,dp
\propto k^{-\alpha}.
$$

\begin{itemize}
    \item $f(p)$: la distribución de la probabilidad de éxito de un solo intento para los distintos problemas;
    \item $\alpha$: el exponente de la cola larga cercana a cero;
    \item $k^{-\alpha}$: la tasa de fallo promedio, que decae conforme crece el presupuesto de muestreo.
\end{itemize}

\videofigure{figures/long-tail-theorem.jpg}{La gran cantidad de problemas con baja probabilidad de éxito hace que el pass@$k$ promedio del dataset siga una ley de potencias.}{00:09:40--00:11:42}

\begin{importantbox}{La cola larga determina la duración del inference scaling}
Si el $p_i$ de todos los problemas es alto, la saturación llega después de pocos muestreos; si una gran cantidad de problemas tiene un $p_i$ extremadamente bajo pero distinto de cero, el beneficio se extiende hasta valores de $k$ muy grandes. Por eso la forma del inference scaling no la determina únicamente el modelo, sino también la distribución de dificultad del benchmark.
\end{importantbox}

\subsection{El cambio en el paradigma de asignación de cómputo}

Cuando el cómputo de inferencia puede intercambiarse de manera predecible por tasa de éxito, el sistema ya no tiene que fijar todo su presupuesto en el pre-training. Para un número reducido de solicitudes de alto valor, invertir una gran cantidad de test-time compute puede resultar más económico que entrenar y desplegar un modelo más grande.

\videofigure{figures/compute-allocation.jpg}{El cambio de "entrenar un modelo grande, inferencia única" hacia "un modelo de tamaño moderado, invirtiendo cómputo de inferencia según la solicitud".}{00:11:42--00:12:38}

\subsection{Resumen de la sección}

El beneficio de Repeated sampling proviene de la cola larga en la probabilidad de éxito de los problemas. Tiene una aplicabilidad amplia respecto de la escala del modelo, pero lo que demuestra es que la oracle coverage es escalable; sin un verifier confiable, un candidato correcto todavía puede no llegar a convertirse en la salida final.

